\documentclass[a4paper, 12pt]{article}
\usepackage[T1]{fontenc} % Für griechische Buchstaben
\usepackage[utf8]{inputenc}
\usepackage[ngerman]{babel}
\usepackage{csquotes}
\usepackage{array}

%Literaturverzeichnis
\usepackage[
    backend=biber, 
    natbib=true,
    style=ext-authoryear, 
    citestyle=apa, 
    autocite=footnote, 
    isbn=false, 
    giveninits=true,
    innamebeforetitle=true,
    maxbibnames=3, 
    maxcitenames=3
]{biblatex}
\addbibresource{UFM89 Assignment.bib} % Muss angepasst werden
% Doppelpunkt nach Jahreszahl
\renewcommand*{\labelnamepunct}{\addcolon\space}    
% Komma statt Punkt zwischen Titel, Ausgabe usw.
\renewcommand*{\newunitpunct}{\addcomma\space}	
% Entfernt Anführungszeichen
\DeclareFieldFormat*{title}{#1}
%Hrsg. in Klammern bei Collections
\DeclareFieldFormat{editortype}{\mkbibparens{#1}}
%Komma vor (Hrsg.) weg
\DeclareDelimFormat{authortypedelim}{\addspace} % Kein Komma nach Autor
\DeclareDelimFormat{editortypedelim}{\addspace}  % Kein Komma nach Hrsg.
% Darstellung der urls in gleicher font
\usepackage{url}
\urlstyle{same}

%Für aktuelles Datum
\usepackage{datetime}
\newdateformat{mydate}{\THEDAY. \monthname[\THEMONTH] \THEYEAR}

% Visuals
\usepackage{setspace} % Für Zeilenabstand
\onehalfspacing % Zeilenabstand 1,5
\usepackage[top=3cm, bottom=3cm, left=2.5cm, right=2cm]{geometry} %Ränder des Dokuments
\setlength{\parindent}{0pt} % space at start of paragraph
\setlength{\parskip}{0.3cm} % space between paragraphs
\usepackage[bottom]{footmisc} % Damit Fußnoten an Seitenende
\usepackage{amsmath}
\usepackage{graphicx} % Zum Einfügen von Bildern
\usepackage{hyperref} %Zum Erstellen von Links
\hypersetup{
    colorlinks=true, 
    allcolors=black
}

\usepackage{lipsum} % Paket für Standardtext
\usepackage{comment}

% Format der Fußnoten ändern
\usepackage{scrextend}
\deffootnote{0.5em}{0em}{\textsuperscript{\thefootnotemark}\,}
\renewcommand{\footnoterule}{%
  \kern -3pt
  \hrule width 10cm 
  \kern 2pt
}

\begin{document}
    \begin{titlepage}
    \begin{center}
        \includegraphics[width=0.4\textwidth]{images/university.png}
        
        \vspace*{4cm}
        
        \Huge
        \textbf{Assignment Template}

        \vspace{0.2cm}
        \large
        Impulse des Neuroleadership-Ansatzes für ein erfolgreiches Change-Management

        \vfill

        Assignment für das Modul\\
        UFM89-Management von Teamwork, Kollaboration und Veränderungsprozessen

        \vspace{0.8cm}
        \normalsize
        \textbf{Fabian Kuhn}\\
        Hauptstraße 16\\
        67133 Maxdorf\\
        \href{mailto:fabian.kuhn1@stud.akad.de}{fabian.kuhn1@stud.akad.de}\\
        Immatrikulationsnummer: 8150162\\
        Bearbeitungszeitraum: 11.12.2024 - 05.02.2025 \\
    \end{center}
\end{titlepage}
    \newpage

    \pagenumbering{Roman}
    
    \tableofcontents
    \newpage

    \input{chapters/0-Abkürzungen.tex}
    \newpage

    \pagenumbering{arabic}
    
    \section{Einleitung}

\subsection{Relevanz und Problemstellung}
Seit seiner Erfindung im Jahre 1976 durch Nikolaus August Otto hat der Viertakt-Ottomotor als Antriebsaggregat eine weitreichende Verwendung gefunden.\footcite[Vgl.][1]{merker2019} 
Obwohl die Elektromobilität an Relevanz gewinnt, setzt sich diese Bedeutung auch heute weiter fort.
Daten des Kraftfahrt-Bundesamtes zeigen, dass Benzinmotoren bei den Neuzulassungen in Deutschland weiterhin einen erheblichen Anteil ausmachen.\footcite[Vgl.][]{kba2025}
Die hohe Energiedichte, sowie einfache Möglichkeiten zum Transport und Speichern machen Flüssigbrennstoff zur bevorzugten Energiequelle für den Verkehr.\footcite[Vgl.][2]{leach2020}
Doch die Anforderungen haben sich angesichts des Klimawandels und schlechter Luftqualität durch $\text{CO}_2$-Emissionen gewandelt.
Eine hohe Last und Motordrehzahl erhöhen die Emission von $\text{CO}_2$\footcite[Vgl.][6]{shahad2015}, weshalb eine präzise Regelung der Motordrehzahl unter variierender Last einen großen Stellenwert einnimmt.

Die experimentelle Entwicklung von Regelungsstrategien am physischen Prüfstand erweisen sich als aufwändiger Prozess.
An dieser Stelle ist die Modellbildung und Simulation ein unverzichtbares Werkzeug.
Dynamische Systemmodelle ermöglichen die virtuelle Nachbildung und Analyse des komplexen Zusammenspiels von Thermodynamik, Mechanik und Regelungstechnik.

Aus diesem Grund widmet sich diese Arbeit der Erweiterung eines bestehenden thermodynamischen Modells eines Viertakt-Ottomotors zu einem dynamischen System.
Ziel wird es sein, das Verhalten des Motors bei einem abrupten Lastwechesel zu simulieren.
Darauf aufbauend soll mittels PID-Reglers eine Drehzahlstabilisierung implementiert werden.
Damit soll demonstriert werden, wie der Motor stets im optimalen Betriebspunkt gehalten wird, um die Energieeffizienz zu steigern.

\subsection{Vorgehensweise und Aufbau der Arbeit}


    
    \section{Grundlagen}

\subsection{Modell nach Scherf}
Das Hauptziel des Modells nach Scherf\footcite[Vgl.][94-101]{scherf2010} ist es den zeitlichen Verlauf von Druck und Temperatur des Gases im Zylinder eines Viertakt-Ottomotors zu simulieren.
Das Modell beschränkt sich dabei auf den Verdichtungs- und Arbeitstakt, da in diesen die wesentlichen thermodynamischen Prozesse stattfinden.
Der Verdichtungstakt beginnt mit dem unteren Totpunkt (UT) bei 180° und endet am oberen Totpunkt (OT) bei 360° Kurbelwinkel.
Der Arbeitstakt beginnt am OT bei 360° und endet am UT bei 540° Kurbelwinkel.
Das Ansaugen und Ausstoßen des Gases wird bewusst vernachlässigt, weshalb die Simualtion bei 540° Kurbelwinkel stoppt.

\subsection{Differenzialgleichung der Temperatur}
Das Kernstück des zugrundeliegenden Modells nach Scherf ist der 1. Hauptsatz der Thermodynamik für ein geschlossenes System. 
Dieser besagt, dass die Änderung der inneren Energie des Gases $\Delta U$ die Summe aus der zugeführten Wärme $\Delta Q$ und der am Gas verrichteten Arbeit $\Delta W$ ist.

In der Simulation wird dies so umgeformt, dass sich eine Differenzialgleichung für die Temperaturänderung $\Delta T / \Delta \varphi$  ergibt. 
Die Gleichung folgt der Idee, dass sich die Temperatur im Zylinder aufgrund von drei Energieflüssen ergibt, welche im Summenblock zusammengeführt werden: der zugeführten Wärme durch die Verbrennung $\Delta \dot{Q}_B$, der abgeführten Wärme an die Zylinderwand $\Delta \dot{Q}_W$, sowie der abgeführten Energie durch die mechanische Arbeit des Kolbens. 
Teilt man diese drei Leistungen durch die Wärmekapazität ergibt sich die Temperaturänderung. 
Durch eine Integration dieser erhält man die absolute Temperatur.

\subsection{Grundlagenthema 3}

    
    \section{Hauptteil}


    \newpage

    \section{Fazit}

\subsection{Zusammenfassung}


\subsection{Ausblick}


\subsection{Limitationen und kritische Reflexion}
Durch den begrenzten Umfang der vorliegenden Arbeit wurden einige Aspekte nur oberflächlich behandelt. 
Die vorgegebenen Ziele konnten dennoch unter Berücksichtigung der Limitationen erreicht werden.
\begin{itemize}
    \item Bei der Literaturrecherche gab es keine Beschränkung auf Zeitschriftenartikel mit Peer-Review. Diese sollten bevorzugt verwendet werden. Zu Grundlagen und Anwendungen konnten nicht genügend relevante Zeitschriftenartikel mit Peer-Review gefunden werden, weshalb auch Monographien und weitere Quellen berücksichtigt wurden.
    \item Die Arbeit basiert ausschließlich auf einer Literaturrecherche, wodurch praktische Erfahrungen fehlen.
    \item Neurowissenschaftliche Erkenntnisse sind oft hochkomplex. Es besteht die Gefahr, dass vereinfachte Modelle wie das SCARF-Modell die psychologischen Prozesse, die bei Veränderungen ablaufen, nur teilweise erfassen.
    \item Soziale Bedürfnisse sind stark subjektiv, was eine universelle Anwendung des Modells schwierig macht.
\end{itemize}

    \newpage

    \pagenumbering{Roman}
    \setcounter{page}{2}

    %Literaturverzeichnis
    \printbibliography
    \addcontentsline{toc}{section}{Literaturverzeichnis}

    \begin{table}[h]
        \centering
        \begin{tabular}{|c|m{15em}|}
            \hline
            KI-basiertes Hilfsmittel & Einsatzform \\ 
            \hline
            ChatGPT & Hilfestellung zur Gliederung der Arbeit, Übersetzungen, Rechtschreibkorrektur \\ 
            \hline
            Elicit & Quellen- und Literaturrecherche \\ 
            \hline
        \end{tabular}
        \caption{Verwendete KI-basierte Hilfsmittel}
    \end{table}
    \newpage

    \input{chapters/eidestattliche erklärung.tex}
\end{document}
